% Computational supplement to the associated paper.
% Build:  latexmk -pdf SUPPLEMENTARY-MATERIAL-v1.0.0.tex
%
% This supplement documents what the reproducibility bundle computes and how its
% outputs are generated. It makes no claim that is not in the paper.
\documentclass[11pt,a4paper]{article}
\usepackage{lmodern}
\usepackage[T1]{fontenc}
\usepackage{amsmath,amssymb}
\usepackage[margin=1in]{geometry}
\usepackage{booktabs}
\usepackage[hidelinks]{hyperref}

\title{Computational supplement\\
\large $1+3$ covariant CMB anisotropies --- analytical reproducibility bundle v1.0.0}
\author{}
\date{}

\begin{document}
\maketitle

\section{Scope}
This bundle is archived at \texttt{doi:10.25375/uct.34069509}, active on publication of
v1.0.0.

This supplement accompanies the v1.0.0 analytical release. The bundle is an
independent reconstruction, from first $1+3$ covariant principles, of the
almost-Friedmann--Lema\^itre results of Gebbie, Dunsby \& Ellis,
\emph{Annals of Physics} \textbf{282}, 321 (2000). It reproduces; it does not
argue. No new effect is claimed.

\section{The acceptance spine}
Appendix F of that paper, pp.~379--380, Eqs.~(F.1)--(F.4), links the covariant
mode recursion to external treatments written in real coefficients. Reproducing
that chain tests the harmonic normalisation without introducing a
representation.\footnote{The comparison applies the $(2\ell+1)^{-1}$ division
that the text preceding (F.4) specifies and the display omits; (F.4) is otherwise
correct as printed. A misprint, recorded in
\texttt{provenance/ANNALS-II-ANTECEDENT-AND-CORRECTIONS-v1.md}.}

Three external treatments are used. Wilson \& Silk is not reproduced here: a
coefficient transcribed through Appendix~F would be the same source as Appendix~F
rather than an independent fourth.

\section{Frames, and what v1.5.0 then is}
The solution is carried with a \textbf{generic} $u^a$. The explicit frames appear
as choices made at the end, not assumed at the start: the Newtonian frame
$\tilde\sigma_{ab}=0$, which is the specialisation Annals II itself makes at
p.~365 in order to remove the shear terms, and the energy frame $q_a=0$. Carrying
the generic threading and specialising late is the difference the $1+3$ covariant
approach buys; a gauge-fixed treatment fixes the frame before the hierarchy is
solved and cannot exhibit it.

This is also what fixes the status of the next release. Because v1.0.0 is the
general case, \textbf{v1.5.0's restricted case is a recovery and a check} --- it
must return v1.0.0's numbers, and that requirement is the regression test rather
than a new result.

Two consequences are visible in the figures. The decomposition of the temperature
source into primary, Doppler and integrated terms (\S7.1.1, Eqs.~(177)--(179)) is
frame-dependent, while the total is not, so it is drawn in two frames. And
footnote~29 on p.~365 asserts that the Newtonian choice also removes the high-$\ell$
truncation difficulty of Appendix~E; that is a checkable claim, and the
convergence figure checks it in both frames.

\section{Conventions}
The normative sheet is carried, generated rather than copied, at
\texttt{provenance/conventions.md}. Equation numbers in this bundle are the
published numbers.

The harmonic normalisation is
$Q_{A_\ell}=(-k_{\rm phys})^{-\ell}\mathrm{D}_{\langle A_\ell\rangle}Q$. Both
candidate conventions are run end to end; only this one reproduces the real,
opposite-sign form of the external
hierarchies.\footnote{Annals~I prints $(-1)^\ell$ at the corresponding step where
its own definition gives $i^\ell$. The covariant multipole $\tau_{A_\ell}$ is
invariant under the choice, so the misprint never produced a wrong result; the
mode coefficients are not invariant, and the two statements are distinct. Both
are asserted as tests.}

\section{Comparison, and its range}
The $C_\ell$ comparison is Sachs--Wolfe only, so its tolerance is bounded by the
acoustic modulation dropped in reaching it. The range is derived, not chosen:
$2\le\ell\le20$ at 5\%, generated by
\texttt{scripts/derive\_ell\_range.py}.\footnote{That derivation uses the Bessel
identity of \S8.3, which requires $[(m/2)!]^2$ in its denominator. As printed it
is exact where $\Gamma(m/2+1)=1$, so the large-scale result (201) is unaffected;
the horizon-scale normalisation (204) and the $D_\ell$ of (205) are. A misprint,
corrected silently and recorded.}

\section{Reproduction}
\begin{verbatim}
python scripts/run_all.py --strict
\end{verbatim}

\section{Figures}
% One subsection per canonical figure, each naming its generating script.
% Pending: see RELEASE-NOTES-v1.0.0.md.

\end{document}
